\section{A simply connected K\"ahler manifold with two torsion primes}
\label{sec:geometry}

Fix an integer $N\geq 704$. We construct a simply connected compact
K\"ahler manifold $X$ of complex dimension $N$ whose integral homology has
prescribed $2$-primary and $3$-primary torsion. The construction has two
steps. First, quotients of products of elliptic curves provide surfaces
with the required torsion groups. Blowing up projective space along these
surfaces, with two projective-line factors added to one center, then
places that torsion in the required degrees.

\subsection{Two quotient surfaces}

Write
\[
 E=\C/(\Z+i\Z),\qquad
 \zeta_p=\exp(2\pi i/p),\qquad
 \Lambda_2=\Z+i\Z,\qquad
 \Lambda_3=\Z+\zeta_3\Z.
\]
For $p\in\{2,3\}$, let $E_p=\C/\Lambda_p$ and define
\[
 \sigma_p:E\times E_p\longrightarrow E\times E_p,
 \qquad
 \sigma_p(z,w)=(z+1/p,\zeta_p w).
\]
Multiplication by $\zeta_p$ preserves $\Lambda_p$, so this map is
well-defined and holomorphic. It has order $p$, and its nonidentity
powers have no fixed points: translation by $k/p$ on $E$ is nontrivial
for $0<k<p$. Thus
\[
 S_p=(E\times E_p)/\langle\sigma_p\rangle
\]
is a smooth compact connected complex surface. These are classical bielliptic surfaces; see \cite{Serrano1991,Blomme2025} for their quotient descriptions and integral topology. We give the specific computations needed here.

\begin{lemma}
\label{lem:torsion-surfaces}
The surfaces $S_2$ and $S_3$ admit holomorphic embeddings into
$\CP^{44}$ and $\CP^{164}$, respectively. Set
\[
 T_2=(\Z/2)^2,\qquad T_3=\Z/3.
\]
Their integral homology is
\begin{equation}
\label{eq:surface-homology}
 H_i(S_p;\Z)\cong
 \begin{cases}
  \Z,&i=0,4,\\
  \Z^2\oplus T_p,&i=1,2,\\
  \Z^2,&i=3,\\
  0,&\text{otherwise}.
 \end{cases}
\end{equation}
In particular, $b(S_p)=8$.
\end{lemma}

\begin{proof}
We first construct the projective embeddings. On an elliptic curve, a
line bundle of degree three has three independent sections. Riemann--Roch
shows that imposing one or two zeros, with multiplicities allowed,
lowers this dimension by one or two. Its complete linear system is
therefore base-point-free and separates points and tangent directions,
giving an embedding in $\CP^2$. Applying this to both curves and then
using the Segre embedding gives a holomorphic embedding
\[
 \iota_p:E\times E_p\longrightarrow\CP^8.
\]

View a point of $\CP^8$ as the projective class of a nonzero linear
polynomial on $(\C^9)^*$. To the orbit of $x\in E\times E_p$ associate
the degree-$p$ polynomial obtained by multiplying representatives of
the points
$\iota_p(x),\iota_p(\sigma_p x),\ldots,
\iota_p(\sigma_p^{p-1}x)$.
Changing representatives changes the product by a nonzero scalar.
Multiplication of linear polynomials is a holomorphic map
\[
 (\CP^8)^p\longrightarrow
 \mathbb P\bigl(\operatorname{Sym}^p\C^9\bigr),
\]
and the product is invariant under cyclically permuting its factors.
Consequently it defines a holomorphic map from $S_p$ to the projective
space on the right. Unique factorization of polynomials shows that this
map separates orbits.

It is also immersive. Indeed, choose local representatives
$L_0,\ldots,L_{p-1}$ of the orbit points and suppose that a tangent
vector makes their product stationary as a projective point. Its
derivatives then satisfy
\[
 \sum_{j=0}^{p-1}\dot L_j\prod_{k\ne j}L_k
   =a\prod_{k=0}^{p-1}L_k
\]
for some scalar $a$. The orbit points are distinct, so the linear
polynomials $L_j$ are pairwise nonproportional. Restricting the identity
to the hyperplane $L_j=0$ shows that $\dot L_j$ vanishes on that
hyperplane: the product of the remaining factors is a nonzero
polynomial there. Hence $\dot L_j$ is a scalar multiple of $L_j$.
Every projective orbit point is stationary, and the immersion
$\iota_p$ forces the original tangent vector to vanish. The quotient
map is a local biholomorphism, so the induced map on $S_p$ is immersive.
A compact injective holomorphic immersion is an embedding. Finally,
\[
 \dim\operatorname{Sym}^2\C^9=45,
 \qquad
 \dim\operatorname{Sym}^3\C^9=165,
\]
which gives the stated ambient dimensions.

For the homology computation, the universal covering $\C^2\to S_p$
has deck group generated by
\[
 a(z,w)=(z+1/p,\zeta_p w),\qquad
 d(z,w)=(z+i,w),\qquad
 t_\lambda(z,w)=(z,w+\lambda)\quad(\lambda\in\Lambda_p).
\]
Here $a^p$ is translation by $1$ in the first coordinate. The elements
$a,d$ commute and generate a copy of $\Z^2$; conjugation by $a$ on
$\Lambda_p$ is multiplication by $\zeta_p$, while conjugation by $d$
is trivial. Thus the deck group is $\Lambda_p\rtimes\Z^2$, and
abelianization gives
\begin{equation}
\label{eq:surface-abelianization}
 H_1(S_p;\Z)
 \cong\Z^2\oplus\Lambda_p/(\zeta_p-1)\Lambda_p.
\end{equation}
For $p=2$, multiplication by $\zeta_p-1$ is $-2I$ on the lattice.
For $p=3$, its matrix in the basis $(1,\zeta_3)$ is
\[
 \begin{pmatrix}-1&-1\\1&-2\end{pmatrix}.
\]
The latter has determinant $3$ and an entry equal to $1$, so its Smith
normal form has diagonal $(1,3)$. These computations identify the
torsion groups in \eqref{eq:surface-abelianization} as $T_p$.

The complex orientation makes $S_p$ an oriented closed real
four-manifold. Integral Poincar\'e duality and the universal coefficient
theorem give
\[
 \Tor H_2(S_p;\Z)\cong\Tor H^2(S_p;\Z)
 \cong\Tor H_1(S_p;\Z)=T_p,
\]
whereas $H_3(S_p;\Z)\cong H^1(S_p;\Z)$ is free of rank two.
The degree-$p$ covering $E\times E_p\to S_p$ gives
$p\chi(S_p)=\chi(E\times E_p)=0$. Since the ranks in degrees
$0,1,3,4$ are $1,2,2,1$, respectively, the rank in degree two is also
two. This proves \eqref{eq:surface-homology}.
\end{proof}

\subsection{Integral homology under blow-up}

We next recall the blow-up formula with integral coefficients. The
coefficient ring matters here: the centers just constructed have
torsion, and we need that torsion to survive as direct summands. The
following standard formula is recorded, in this integral generality,
in \cite[Theorem~4.1]{DuanLi2009}; we include its proof.

\begin{proposition}
\label{prop:blowup}
Let $A$ be a compact connected K\"ahler manifold, and let $Z\subset A$
be a closed connected complex submanifold of complex codimension
$c\geq2$. Its complex blow-up $\pi:\widehat A\to A$ is compact,
connected, and K\"ahler. If $A$ is simply connected, then so is
$\widehat A$. For every integer $i$ there is an isomorphism of abelian
groups
\begin{equation}
\label{eq:blowup-homology}
 H_i(\widehat A;\Z)
 \cong H_i(A;\Z)\oplus
 \bigoplus_{j=1}^{c-1}H_{i-2j}(Z;\Z).
\end{equation}
\end{proposition}

\begin{proof}
Let $\nu$ denote the complex normal bundle of $Z$ in $A$. In local
coordinates transverse to $Z$, the blow-up replaces a vector in
$\C^c$ by a pair consisting of a line and a vector on that line.
These local models give a smooth complex manifold and a proper
holomorphic map $\pi$, which is an isomorphism away from $Z$. The
exceptional divisor is $E=\mathbb P(\nu)$, with projection
$\rho:E\to Z$, and its normal line bundle in $\widehat A$ is the
tautological bundle $\mathcal O_E(-1)$. The model also shows that
$\widehat A$ is connected; properness over compact $A$ gives
compactness.

To obtain a K\"ahler form, let $\omega_A$ be one on $A$. Choose a
Hermitian metric on $\nu$ and its induced metric on
$\mathcal O_E(1)$, whose curvature is positive on each projective
fiber. Since
$\mathcal O(-E)|_E\cong\mathcal O_E(1)$, this metric extends to a
Hermitian metric on $\mathcal O(-E)$. Let $\eta$ be its real Chern
curvature form. The form $\pi^*\omega_A$ is semipositive; along $E$
its kernel is precisely the tangent space to the projective fiber,
and it is positive away from $E$. On those kernel directions $\eta$
is positive. Compactness of the unit tangent bundle now shows that
\[
 \pi^*\omega_A+\varepsilon\eta
\]
is positive for every sufficiently small $\varepsilon>0$. More
explicitly, $\eta$ is positive in a neighborhood of the kernel
directions, while on the complement $\pi^*\omega_A$ has a positive
lower bound and $\eta$ is bounded. The displayed form is closed and
of type $(1,1)$, hence K\"ahler.

Suppose that $A$ is simply connected. Since $Z$ has real codimension
at least four, general position allows paths and contracting disks
to avoid $Z$; consequently $A\setminus Z$ is simply connected.
Every loop in $\widehat A$ can be perturbed off the real codimension
two submanifold $E$. It then lies in
$\widehat A\setminus E\cong A\setminus Z$ and contracts there.
This proves simple connectivity of $\widehat A$.

It remains to establish the integral homology formula. We first work
in cohomology, always with coefficients in $\Z$. Write
$\iota:E\hookrightarrow\widehat A$ for the inclusion and put
$h=c_1(\mathcal O_E(1))$. By Leray--Hirsch, the map
\[
 \bigoplus_{j=0}^{c-1}H^{k-2j}(Z;\Z)
 \longrightarrow H^k(E;\Z),\qquad
 (u_0,\ldots,u_{c-1})\longmapsto
 \sum_{j=0}^{c-1}h^j\rho^*u_j
\]
is an isomorphism. It uses an integral basis on every fiber and
requires no restriction on torsion in $H^*(Z;\Z)$.

The Gysin map for the oriented normal line of $E$ satisfies
\[
 \iota^*\iota_*(v)=c_1(\mathcal O_E(-1))v=-hv.
\]
For $1\leq j\leq c-1$ define
\[
 s_j:H^{k-2j}(Z;\Z)\longrightarrow H^k(\widehat A;\Z),
 \qquad s_j(u)=-\iota_*\bigl(h^{j-1}\rho^*u\bigr).
\]
Then $\iota^*s_j(u)=h^j\rho^*u$, so these maps lift all the
positive powers of $h$ in the preceding decomposition.

We claim that every class $v\in H^k(\widehat A;\Z)$ with
$\iota^*v\in\rho^*H^k(Z;\Z)$ lies in the image of $\pi^*$.
The map of pairs
\[
 (\widehat A,E)\longrightarrow(A,Z)
\]
induces an isomorphism on relative cohomology: after collapsing the
subspaces, it induces a homeomorphism
$\widehat A/E\cong A/Z$. Write $\iota^*v=\rho^*z$.
Naturality of the connecting homomorphisms of the two pairs implies
that the connecting image of $z$ vanishes, because its pullback is
the connecting image of $\iota^*v$, which is zero. Hence $z$ extends
to a class $u\in H^k(A;\Z)$. The difference $v-\pi^*u$ restricts
to zero on $E$ and therefore comes from $H^k(\widehat A,E;\Z)$.
Pulling that relative class back from $H^k(A,Z;\Z)$ proves the
claim.

Subtracting suitable classes $s_j(u_j)$ from an arbitrary $v$ leaves
a class covered by this claim. Thus
\[
 \pi^*+\sum_{j=1}^{c-1}s_j:
 H^k(A;\Z)\oplus\bigoplus_{j=1}^{c-1}H^{k-2j}(Z;\Z)
 \longrightarrow H^k(\widehat A;\Z)
\]
is surjective. It is injective as well. First, $\pi^*$ is injective
over $\Z$: the map $\pi$ has degree one, so the pushforward defined
using integral Poincar\'e duality is a left inverse to $\pi^*$.
If a sum in the displayed map is zero, restriction to $E$ and
Leray--Hirsch force all its $u_j$ for $j\geq1$ to vanish; injectivity
of $\pi^*$ then forces the remaining class to vanish. We have proved
\begin{equation}
\label{eq:blowup-cohomology}
 H^k(\widehat A;\Z)
 \cong H^k(A;\Z)\oplus
 \bigoplus_{j=1}^{c-1}H^{k-2j}(Z;\Z).
\end{equation}

If $\dim_{\C}A=m$, integral Poincar\'e duality applied to
\eqref{eq:blowup-cohomology} in degree $2m-i$ gives center summands
$H_{i-2c+2j}(Z;\Z)$. Replacing $j$ by $c-j$ yields
\eqref{eq:blowup-homology}.
\end{proof}

\subsection{The manifold \texorpdfstring{$X$}{X}}

The blow-up formula now allows us to place both torsion primes in one
simply connected manifold. The different dimensions of the two centers
will give their torsion contributions different behavior near the ends
of the degree range.

\begin{proposition}
\label{prop:torsion-space}
For the fixed integer $N\geq704$, there is a simply connected compact
K\"ahler manifold $X$ of complex dimension $N$ obtained by blowing up
$\CP^N$ along disjoint copies of
\[
 Z_2=S_2,\qquad Z_3=S_3\times\CP^1\times\CP^1.
\]
Set $L=N-3$. Its only nonzero integral homology torsion groups occur
in degrees $1+2j$ and $2+2j$, where $1\leq j\leq L$, and are
\begin{equation}
\label{eq:X-torsion}
 \Tor H_{1+2j}(X;\Z)\cong\Tor H_{2+2j}(X;\Z)
 \cong(\Z/2)^{a_j}\oplus(\Z/3)^{q_j},
\end{equation}
where
\begin{equation}
\label{eq:torsion-profiles}
 a_j=2,\qquad
 q_j=
 \begin{cases}
  1,&j=1,L,\\
  3,&j=2,L-1,\\
  4,&3\leq j\leq L-2.
 \end{cases}
\end{equation}
Moreover,
\begin{equation}
\label{eq:rank-X}
 b(X)=41N-183.
\end{equation}
\end{proposition}

\begin{proof}
By Lemma~\ref{lem:torsion-surfaces}, $Z_2$ embeds in $\CP^{44}$.
The embedding $S_3\hookrightarrow\CP^{164}$ and the Segre embedding
give
\[
 Z_3\hookrightarrow
 \CP^{164}\times\CP^1\times\CP^1
 \hookrightarrow\CP^{659},
\]
since $165\cdot2\cdot2=660$. The direct-sum decomposition
$\C^{705}=\C^{45}\oplus\C^{660}$ places $\CP^{44}$ and
$\CP^{659}$ in disjoint linear subspaces of $\CP^{704}$.
Including that space linearly in $\CP^N$ gives the required disjoint
centers. Define $X$ by blowing up $Z_2$ and then $Z_3$; the latter
center is unchanged by the first blow-up because they are disjoint.
Both centers have complex codimension at least two, so
Proposition~\ref{prop:blowup} shows that $X$ is compact, connected,
simply connected, and K\"ahler.

The complex codimensions of $Z_2$ and $Z_3$ are $N-2$ and $N-4$,
respectively. The integral blow-up formula therefore gives
\begin{equation}
\label{eq:X-homology}
\begin{split}
 H_i(X;\Z)\cong{}&H_i(\CP^N;\Z)
 \oplus\bigoplus_{r=1}^{N-3}H_{i-2r}(S_2;\Z)\\
 &\oplus\bigoplus_{r=1}^{N-5}
 H_{i-2r}(S_3\times\CP^1\times\CP^1;\Z).
\end{split}
\end{equation}
The first center contributes $(\Z/2)^2$ in each of degrees
$1+2j$ and $2+2j$, for $1\leq j\leq L$. For the second center,
the homology of $\CP^1\times\CP^1$ is free, with ranks $(1,2,1)$
in degrees $(0,2,4)$. K\"unneth therefore has no additional
$\operatorname{Tor}$ terms. In either of the two degree lists, its
$3$-primary multiplicity is
\[
 q_j=
 \sum_{r=1}^{L-2}
 \bigl(\mathbf 1_{j=r}+2\mathbf 1_{j=r+1}
                  +\mathbf 1_{j=r+2}\bigr),
\]
where $\mathbf 1$ denotes the indicator of the stated equality.
At the endpoints this gives $q_1=q_L=1$, next to them it gives
$q_2=q_{L-1}=3$, and at every remaining index it gives $q_j=4$.
There is no torsion outside these lists by
\eqref{eq:surface-homology} and \eqref{eq:X-homology}. Thus the
profiles are
\[
 (a_1,\ldots,a_L)=(2,\ldots,2),\qquad
 (q_1,\ldots,q_L)=(1,3,4,\ldots,4,3,1).
\]
The projective-line factors in the second center produce the multiplicities
$(1,2,1)$: the resulting $3$-primary multiplicity is smaller than the
$2$-primary multiplicity at the two endpoints and larger at every interior
index. This reversal is what the product comparison in Section~\ref{sec:counts}
will exploit.
Finally, $b(S_2)=b(S_3)=8$ and
$b(S_3\times\CP^1\times\CP^1)=32$. Summing ranks in
\eqref{eq:X-homology} yields
\[
 b(X)=(N+1)+8(N-3)+32(N-5)=41N-183.
\]
\end{proof}
